\begin{abstract}
Multi-agent systems built on large language models are written the way
parallel programs were written before 1994: every framework bundles its own
orchestration policy, its own communication mechanism, and its own failure
handling, and none of them compose. The agent-interoperability protocols
that have emerged since 2024---MCP, A2A, ACP---standardise transport,
discovery and tool invocation, but they address a single client talking to a
single server. None of them provides what a \emph{parallel program} needs:
group membership, collective operations, a consistency model for shared
state, or defined semantics when a participant dies.

We argue that this is precisely the gap MPI filled for distributed-memory
computing, and that MPI's answers largely transfer. We present \ampi{}, a
message-passing protocol for building multi-agent harnesses. \ampi{} adopts
MPI's load-bearing abstractions---communicators with isolated contexts,
groups, typed envelopes, the full collective catalogue, one-sided windows,
collective I/O, a name-shifted profiling interface, and ULFM-style
revoke/shrink/agree---and changes exactly what the executor forces us to
change.

Four properties of agents break MPI's assumptions, and each drives a
specific protocol change. Context is a \emph{cumulative} resource rather
than a reusable buffer, so datatypes carry token bounds, receives are
subject to admission control, and collective algorithm selection must reason
about feasibility before cost; we show that a capacity-aware reduction tree
must be planned against the root's cumulative ingest rather than its
per-round ingest, a distinction with no MPI counterpart. Reduction operators
are semantic rather than arithmetic, so operators declare commutativity,
idempotency and an output bound, and the runtime refuses schedules those
declarations do not support---we show that the natural choice of a
commutative merge operator for a shared-terminology scan is
\emph{unsound}, because a commutative operator cannot express precedence.
Failure is not fail-stop but a lattice of crashed, stalled, drifted and
bankrupt, so the failure detector separates liveness from progress and
recovery prefers rank replacement over shrinking. And because computation
costs five to seven orders of magnitude more than communication, the
latency--bandwidth trade-offs that shape MPI's algorithm selection invert:
logarithmic schedules win everywhere.

We implement \ampi{} as a transport-independent runtime with a shared-directory
device and a command-line binding, so that any process able to run a shell
command---including a coding agent---is a rank. We evaluate it with genuine
multi-agent runs: a twelve-rank book translation in which an exclusive
parallel-prefix scan propagates terminology in $\lceil\log_2 p\rceil$ rounds
instead of $p-1$, an eight-module collaborative software build judged by a
frozen integration suite, and fault-injection runs in which ranks are killed
mid-collective. Microbenchmarks scale the protocol to 128 ranks.
\end{abstract}
